\chapter{A Roman Quest}
"If one might only have a guide to the truth."--\textit{Seneca}.


On Scopus, the high mountain north of Jerusalem, the Roman camp was
pitched, that last autumn in the ministry of Jesus of Nazareth.  A
few years further on, if the warriors of the Emperor Tiberius could
then have foreseen the future, Titus was to quarter his famous
legions on that vantage point; and from its elevation he was to
hurl himself as a resistless battering ram against the Holy City.
But, on this autumn day, when these chronicles begin, no blare of
trumpets was summoning the Roman soldiery to arms; only the feet of
the camp sentinels, as they walked their appointed rounds, broke
the quiet of the sunlit afternoon.

That lithesome, cultivated, serious-minded young knight, Quintus
Cornelius Benignus, is standing on the height which overlooks the
great metropolis.  He is the son of Marcus Cornelius Magnus, that
Roman noble who is the intimate associate of the reigning Caesar,
and who has been a luxurious resident on the Palatine Hill since
his distinguished proconsulship in Africa.

     *     *     *     *     *

NOTE.--It is not from any time-marked Hebrew roll that this story
of Quintus is now taken.  He was of Roman blood, and his record is,
rather, to be found in the Latin literature of his time.  Well it
is when some new leaf is discovered among the musty folios,
reciting the saintly character and the triumphs of those who lived
when Christianity was new.  This record shows the worth of
consecrated life and service in the days when the luxurious Roman
state most needed a Christian citizenship.  But the lesson is none
the less for these last days, when the hope of the world is in the
creed of Quintus.

     *     *     *     *     *

By the side of Quintus is his fellow soldier Aulus.  They had spent
their boyhood together among the scenes of Rome; now they are
companions still, on this last Roman expedition to the district of
Judaea.  While the common soldiery are throwing their  dice  in
the  camp thoroughfare, these are speaking of more serious things.
The picture on which they look from lofty Scopus includes the
shining roofs of Jerusalem, the wooded Mount of Olives, and the far
landscape to the south and west; its undulations and brilliant
colorings no Roman artist might put upon the canvas.

With the autumn haze covering the extended panorama, Quintus says
first to his comrade:

"What the fates have in store for me, here in the city of
Hierosolyma, I am much wondering.  The day before our trireme
sailed from Brundisium for Tyrus I made a visit to the augur's
tent.  His prediction was that my journey hither would be followed
by strange consequences.  The flight of the birds through the air
did not reveal to him just what was to occur; but that something
eventful was to take place he was very sure.  What is to be my
fortune?"

"Your lot it may be," answers Aulus, "to perform some daring deed,
here in our Jewish campaign; and on your return to Rome you may
receive a great reward from the hand of Tiberius."

"In my mind this has been," replies Quintus; "before I left Rome I
had an audience with our divine Caesar, and he was pleased to say
that my fidelity here might bring me special recompense.  Yet would
that be satisfying?  I have seen the triumphal processions in the
streets of Rome, when heroes have been acclaimed; I have heard our
statesmen in the Senate hall, and prize the joys of oratory; I have
been served all my days by slaves in my father's palace, and know
the sweetness of the Falernian wine in the banquet room.  A
proconsulate, if I might come to that dignity, would be a high
honor to write in my life story.  But, my dear Aulus, would there
be content in this?  My restless soul seems crying out for some
better gift from the gods."

"It cannot be," continues Aulus.  "that your heart's love is
involved.  When our military movements bring the Roman knights to
Palaestina, in their pride of birth they do not wed the black-eyed
daughters of the Jews.  On your earlier expedition to Egypt you met
a princess of the land, but were not let to espouse that swarthy
maiden of the Nile.  The reward of love cannot be the experience of
which the augur spoke at Brundisium."

"Not so," says Quintus in response; "as I was leaving Rome, it was
the beautiful Lucretia who sent me forth with her rare farewell.
For my return from Palaestina she is now waiting; and under the
blue skies of Italia we are to wed.  I have been wondering,"
Quintus adds further, "if the augur, watching the flight of birds
there at Brundisium.  could mean that I am to fall by death, here
in Palaestina.  We have not come for battle, but to guard the
peace.  Yet it is easy for Atropos, that cruel fate, to clip the
slender thread of life and send men on to die land of shades.  If
this was what the augur meant, no Roman in the days of Tiberius has
ever set forth upon a more serious adventure."

"You are given to melancholy, this autumn afternoon, my comrade
Quintus," the other says; "you are feeling that sadness which comes
to men when the Dryads move over the earth and touch the leaves
into crimson and gold and brown."

"Not so," answers Quintus; "but I am remembering that I have come
into a land where a strange Teacher is speaking to men of a future
life.  Yet are men to live again?  I have seen the marble tombs on
the Appia Via where the Scipios, the Metelli, and so many more of
our great Romans lie asleep.  Shall I soon follow them?  Is it an
endless slumber?  What is it that the new Rabbi from Nazareth
means, when in the city yonder he speaks of another life?"

"A fig for your weird autumn fancy," responds Aulus; "down to the
streets of Hierosolyma we will go, and among their novel sights we
will forget your serious meditations."


They walk that afternoon as sightseers through the crowded Jewish
emporium.  The shops remind them, with all their contrasts, of the
marts of Rome, for men always and everywhere have the trader's
passion.  In the narrow streets of Jerusalem they see the stir of
many activities.  The workman is hammering his brass; the shoemaker
shapes his sandals; the flax spinner is winding his thread; the
scribe sits on his mat, and is ready for his writing.  In the shops
they see costly merchandise for sale--silks and jewels, fine linens
and perfumes, delicious foods and drinks.  These have been imported
from far Arabia and India; they have been brought from distant
Persia and Media.  With all their variety, no taste, however
fitful, need go unsatisfied.

What a motley crowd is on the streets!  They hear the Aramaic
speech of Palestine, which Quintus has been taught by his Athenian
tutor, and their ears also catch the accents  of other  foreign
tongues.  They meet traders from western Zidon, sailors from Crete,
bearded Idumaeans from beyond Judaea, and scholars from far
Alexandria.  Magnificent Jerusalem it is!  Yet destined soon to
fall.  For the day draws near when the Roman Titus shall weep on
Scopus over its fading splendors and then shall smite it to the
dust.

One purchase only does Quintus make.  In a shop where Egyptian
wares are sold he says to Aulus:

"Look on this scarab, this sacred beetle, which has been shaped by
some workman down in Thebae on the Nile.  We may be sure that no
people believes more intensely in a future life.  What compliment
they pay this physical frame of men when they hold that embalmment
restores to the soul its former body!  After the judgment of
Osiris, if their lives be true, the worthy shall enjoy the
companionship of the great god forever.  No other people wears such
a visible emblem of their faith in another life.  I will buy this
scarab for an amulet against accident and evil."


But where had the workman gone who once had shaped that token of
immortality?  Whither had vanished his carver's skill?  Where had
disappeared his projects and his dreams?  Quintus is not thinking
of any proconsulship he may win, or even of the love light in the
eyes of Lucretia, as he climbs again the heights of Scopus.  Rather
he is meditating on the departed maker of scarabs--and on the
destiny of the soul.  For ages the philosophers have been
speculating about the future life.  Familiar is Quintus with the
views of Laelius and Seneca, among the Roman inquirers, and with
the teachings of the great Grecians who have spoken in classic
Athens.  But now the question leaps to the front.  Quintus is in
the city where Ayran travelers and Persian magi and Egyptian
priests are busy telling their theories of immortality.  He is in
the very streets, besides, where a sandaled Teacher from Nazareth
is declaring that the dead shall live again.  If but half is true
that this strange Man is reputed to have said, no priest of Jupiter
has ever uttered at Rome so luminous a word.  Can it be that
Quintus himself shall see this Christus and hear his message?  If
so, his will be in very truth a momentous quest.